\documentclass[10pt]{article}
\usepackage[margin=0.61in]{geometry}
\usepackage{amsmath,amsthm,amssymb,microtype}
\usepackage{titlesec}
\usepackage{fontspec,unicode-math}
\setmainfont{TeX Gyre Pagella}
\setsansfont{TeX Gyre Heros}
\setmonofont{Latin Modern Mono}
\setmathfont{TeX Gyre Pagella Math}
\pdfvariable suppressoptionalinfo 512\relax
\pdfvariable trailerid {[<c338c338c338c338c338c338c338c338><c338c338c338c338c338c338c338c338>]}
\usepackage[hidelinks]{hyperref}
\newtheorem{theorem}{Theorem}
\newtheorem{lemma}{Lemma}
\providecommand{\CRevisionRound}{2}
\newcommand{\Prb}{\mathbb P}
\newcommand{\E}{\mathcal E}
\newcommand{\T}{\mathcal T}
\newcommand{\LE}{\operatorname{LE}}
\newcommand{\diag}{\operatorname{diag}}
\titlespacing*{\section}{0pt}{1.1ex plus .2ex}{0.4ex}
\titlespacing*{\subsection}{0pt}{0.8ex plus .2ex}{0.3ex}
\titlespacing*{\paragraph}{0pt}{0.45ex}{0.5em}
\title{From Abelian Cycle Popping to Every Transfer-Current Minor}
\author{Anonymous}
\date{3 September 2026}
\begin{document}
\maketitle
\vspace{-1.6em}

\begin{abstract}
On a finite connected loopless conductance multigraph, we prove one
convention-complete theorem linking infinite-stack cycle popping, Wilson's
loop-erased random walks, weighted spanning trees, the matrix-tree
normalization, and all transfer-current edge determinants.  A local diamond
and strip argument makes every legal cycle rule terminate with identical pop
counts once one terminating order exists; Wilson exploration supplies such an
order almost surely.  A labelled last-exit formula then telescopes killed
Green-function cofactors to give the exact weighted tree law.  Finally,
multivariate conductance perturbations yield every joint edge-inclusion minor,
not only effective-resistance marginals.  The statement includes the
singleton, already-a-tree, distinctly labelled parallel-edge, orientation,
and root-change boundaries.  Exact evidence checks all 772 connected labelled
simple graphs through five vertices, all 8,136 graph--tree pairs and 55,895
edge-subset events, plus 24 weighted multigraphs and 12,754 finite stack
tables.  Enumeration is only a reproducibility receipt: the proof owns the
infinite-stack result.  No arithmetic or target-zeta interpretation is made.
\end{abstract}

\paragraph{Revision certificate.}
\ifcase\CRevisionRound
Core stochastic closure: abelian stacks, Wilson paths, and the weighted law.
\or
Determinantal boundary closure: all transfer-current minors and edge conventions.
\else
Route firewall receipt: exhaustive exact evidence, ownership, and nonclaims.
\fi

\section{Model and theorem}
Let $G=(V,E)$ be a finite connected loopless undirected multigraph.  Parallel
edges remain distinctly labelled and have conductances $c_e>0$.  Freeze one
orientation per edge,
\[
 b_e=\delta_{\mathrm{tail}(e)}-\delta_{\mathrm{head}(e)},\qquad
 B=(b_e)_{e\in E},\quad C=\diag(c_e),\quad L=BCB^{\mathsf T}.
\]
Fix a root $r$ and put $c(v)=\sum_{e\ni v}c_e$.  At each $v\ne r$ place an
independent infinite stack $(X_{v,j})_{j\ge1}$ with
$\Prb(X_{v,j}=e)=c_e/c(v)$ for $e\ni v$.  Its visible card points from $v$
across $e$.  A legal update selects a visible directed cycle and advances
exactly the stack pointers on that cycle.

Let $L^+$ be the Moore--Penrose pseudoinverse and define
\[
 H_{ef}=c_f b_e^{\mathsf T}L^+b_f,\qquad
 K_{ef}=\sqrt{c_ec_f}\,b_e^{\mathsf T}L^+b_f.                 \tag{1}
\]
$H$ is convenient over rational conductances; $K$ is symmetric.  They are
diagonally similar and have the same principal minors.

\begin{theorem}\label{thm:main}
For every fixed stack realization for which one finite legal pop sequence is
terminal, all legal cycle choices terminate with the same vertexwise pop
counts and terminal arborescence.  Under the independent stack law this
happens almost surely, and a canonical legal order is Wilson chronological
loop erasure for any ordering of $V\smallsetminus\{r\}$.

The unoriented output satisfies, for every labelled spanning tree $T$,
\[
 \Prb(\T=T)=\frac{\prod_{e\in T}c_e}{Z_c},\qquad
 Z_c:=\sum_{S\text{ tree}}\prod_{e\in S}c_e=\det L^{(r)},     \tag{2}
\]
where $L^{(r)}$ deletes the root row and column.  For every set of pairwise
distinct labelled edges $S=\{e_1,\ldots,e_k\}$,
\[
 \Prb(S\subseteq\T)=\det(H_{e_i e_j})_{i,j=1}^k
                    =\det(K_{e_i e_j})_{i,j=1}^k.            \tag{3}
\]
Equations (2)--(3) are independent of the root and frozen orientations.
\end{theorem}

\section{Abelian stacks and Wilson equivalence}
\begin{lemma}[diamond and strip]\label{lem:abelian}
If a stack configuration admits a finite terminal legal sequence, then every
legal sequence is finite and has the same pop-count vector and terminal state.
\end{lemma}
\begin{proof}
Two visible directed cycles that share a vertex coincide, since every nonroot
vertex has one visible outgoing arrow.  Distinct visible cycles are therefore
vertex-disjoint; popping either leaves the other visible, and the two pointer
increments commute.

Fix a terminal list $\alpha=(C_1,\ldots,C_m)$ and any cycle $D$ visible at the
initial state.  Unless $D=C_1$, the cycles are disjoint, so $D$ commutes past
$C_1$ and remains visible.  Continue along $\alpha$.  It must eventually equal
one of its cycles; otherwise $D$ would remain visible after the terminal list.
Thus $D$ can be moved to the front of a terminal list without changing any
pointer increment.  Induction on $m$, after each arbitrary first choice,
shows that every legal list has the same finite multiset of pointer increments
and terminal state.  An acyclic functional digraph with only $r$ lacking an
outgoing arrow is an arborescence toward $r$: every forward orbit ends at $r$,
or else repeats and makes a cycle.
\end{proof}

Start with the built tree $\{r\}$.  From the first ordered vertex outside it,
follow visible arrows.  When the path revisits a vertex, its intervening
arrows are a visible cycle; pop them and retain the chronological loop
erasure.  Upon first hitting the built tree, adjoin the retained path and
continue.  Newly exposed cards are fresh conductance transitions, so each
exploration is the random walk on the finite irreducible graph stopped on a
nonempty set.  It hits that set in finite time almost surely.  Finitely many
explorations give a terminal legal list, card for card equal to Wilson's
algorithm.  Lemma~\ref{lem:abelian} promotes this one almost surely finite
order to every legal rule.

\section{Last exits, telescoping, and the weighted law}
Let
\[
 P(x,y)=\sum_{e:x\leftrightarrow y}\frac{c_e}{c(x)}.
\]
For a target set $A$, take a labelled self-avoiding path
$\gamma=(x_0,e_0,x_1,\ldots,e_{q-1},x_q)$ with $x_q\in A$ and earlier
vertices outside $A$.  Set
$D_i=V\smallsetminus(A\cup\{x_0,\ldots,x_{i-1}\})$ and
$G_D=(I-P_D)^{-1}$.

\begin{lemma}[labelled last exits]\label{lem:last}
\[
 \Prb(\LE=\gamma)=\prod_{i=0}^{q-1}
 G_{D_i}(x_i,x_i)\frac{c_{e_i}}{c(x_i)}.                     \tag{4}
\]
\end{lemma}
\begin{proof}
Before the retained step from $x_i$, sum over all excursions that start and
return to $x_i$ without leaving $D_i$.  Their total mass is the diagonal
killed Green entry.  The next retained labelled edge has probability
$c_{e_i}/c(x_i)$.  The strong Markov property after that last exit repeats the
argument in $D_{i+1}$, proving (4).
\end{proof}

Cramer's rule gives
\[
G_D(x,x)=\frac{\det(I-P_{D\smallsetminus\{x\}})}{\det(I-P_D)}. \tag{5}
\]
The factors (5) telescope along each Wilson path and then across all phases,
from $D=V\smallsetminus\{r\}$ to the empty set.  A fixed tree oriented toward $r$
uses one labelled edge $e(v)$ from each $v\ne r$, hence
\[
 \Prb(\T=T)=
 \frac{\prod_{v\ne r}c_{e(v)}/c(v)}{\det(I-P^{(r)})}
 =\frac{\prod_{e\in T}c_e}{\det L^{(r)}},                  \tag{6}
\]
because $I-P^{(r)}=D_r^{-1}L^{(r)}$ for
$D_r=\diag(c(v):v\ne r)$.  The output is almost surely exactly one tree.
Summing (6) over all labelled trees proves both the normalization and the
matrix-tree identity in (2); neither was inserted as an assumption.

\ifnum\CRevisionRound>0
\section{All transfer-current minors}
Delete the root row of $B$ to obtain $R$, and put
$A=RCR^{\mathsf T}=L^{(r)}$.  If $b$ has coordinate sum zero, solutions of
$Lx=b$ differ only by constants.  Incidence vectors annihilate constants, so
\[
 b_e^{\mathsf T}L^+b_f=R_e^{\mathsf T}A^{-1}R_f,
 \qquad H=R^{\mathsf T}A^{-1}RC.                            \tag{7}
\]
This also proves that (7) is independent of the deleted root.

Introduce one commuting indeterminate $t_e$ per labelled edge and replace
$c_e$ by $c_e(1+t_e)$.  From (2), the normalized partition polynomial is
\[
 \frac{Z(c_e(1+t_e))}{Z(c_e)}
 =\sum_T\Prb(\T=T)\prod_{e\in T}(1+t_e).                   \tag{8}
\]
The matrix determinant lemma and Sylvester's determinant identity give the
second representation
\[
 \frac{\det(A+RC\diag(t)R^{\mathsf T})}{\det A}
 =\det(I+\diag(t)R^{\mathsf T}A^{-1}RC)
 =\det(I+\diag(t)H).                                       \tag{9}
\]
The coefficient of $\prod_{e\in S}t_e$ in (8) is
$\Prb(S\subseteq\T)$; its coefficient in (9) is the principal minor
$\det H_S$.  This proves (3) for every $k$ at once.  If
$D=\diag(\sqrt{c_e})$, then $K=DHD^{-1}$, while reversing frozen edge
orientations conjugates either kernel by a diagonal sign matrix.  Principal
determinants are unchanged.

\paragraph{Boundary closure.}
For a singleton, $E=\varnothing$, the empty tree and empty determinant both
have value one.  If $G$ is already a tree, it is the only output and the
rank-$|E|$ projection $K$ is the identity.  Parallel labels retain separate
conductances; their incidence columns agree up to sign, so a minor containing
two parallel edges vanishes, exactly as simultaneous tree inclusion does.
Changing the root changes stack directions but not (2), (7), or any
unoriented event.  Loops and nonpositive conductances lie outside the model.
\fi

\ifnum\CRevisionRound>1
\section{Exact receipt, ownership, and scope}
The exact ledger covers all 772 connected labelled simple graphs on at most
five vertices, all 8,136 graph--tree pairs, and all 55,895 edge-subset events.
It adds 24 positive-integer weighted multigraphs containing labelled parallel
edges: 846 weighted trees and 7,032 further subset events.  A separate
all-schedule audit explores 12,754 depth-two stack tables over every root of
every connected labelled simple graph through four vertices; every terminal
table has one pop-count vector and agrees with canonical Wilson exploration.
Producer-independent reconstruction, symbolic triangle/parallel/$K_4$
identities, isolated byte replay, repaired-hash mutation, and optimized-Python
refusal close the executable boundary.  These finite checks audit conventions,
not almost-sure termination.

Wilson owns the LERW/cycle-popping lineage; Burton--Pemantle own the
transfer-impedance theorem lineage; Kirchhoff and Chaiken own the determinant
lineage.  This source-local reconstruction claims no priority.  C176 concerns
sandpile translations, and C181 deterministic rotor-router torsors; neither
owns independent resampling stacks plus the weighted UST edge process.

The Route-A tuple is $(\mathrm{A0\_FAIL},\mathrm{A1\_FAIL},
\mathrm{A2\_FAIL},\mathrm{A3\_FAIL},\mathrm{A4\_FORMAL\_HINT})$, overall
rejected.  The finite projection in (1) is only a formal source hint.  The
tree partition polynomial is not a target Euler factor or target zeta; no
target local data, root number, automorphy, divisor, functional equation,
zero match, RH claim, Hilbert--P\'olya operator, or Route B invocation occurs.
\fi

\paragraph{Verified source lineage.}
\scriptsize\raggedright
D. B. Wilson, ``Generating random spanning trees more quickly than the cover
time,'' STOC (1996), 296--303,
\href{https://doi.org/10.1145/237814.237880}{doi:10.1145/237814.237880}.
R. Burton and R. Pemantle, ``Local Characteristics, Entropy and Limit Theorems
for Spanning Trees and Domino Tilings Via Transfer-Impedances,'' \emph{Ann.
Probab.} 21 (1993), 1329--1371,
\href{https://doi.org/10.1214/aop/1176989121}{doi:10.1214/aop/1176989121}.
G. Kirchhoff, \emph{Ann. Phys.} 148 (1847), 497--508,
\href{https://doi.org/10.1002/andp.18471481202}{doi:10.1002/andp.18471481202}.
S. Chaiken, \emph{SIAM J. Algebraic Discrete Methods} 3 (1982), 319--329,
\href{https://doi.org/10.1137/0603033}{doi:10.1137/0603033}.
R. Lyons and Y. Peres, \emph{Probability on Trees and Networks}, CUP (2016),
\href{https://rdlyons.pages.iu.edu/prbtree/}{official author page}.

\end{document}
